\section{Approximating the distribution at a fixed time}\label{sec: Jain-Sawhney}

The goal of this section is to prove the following theorem, which provides a detailed reformulation of \Cref{thm: X and nu}.

\begin{thm}\label{thm: restate of X and nu}
Let $n_1,n_2,\ldots$,$k_1,k_2,\ldots$, and $t_1,t_2,\ldots$ be sequences of positive integers such that:
\begin{enumerate}
\item $\lim_{N\rightarrow\infty}n_N=\infty$.
\item $2\le k_N\le n_N$ for all $N\ge 1$, and $\lim_{N\rightarrow\infty}\frac{k_N}{n_N/(\log n_N)^4}=0$.
\item $|t_N-\frac{n_N\log n_N}{k_N}|\le \frac{n_N}{2k_N}\log\log\log n_N$ for all $N\ge 1$. 
\end{enumerate}
Define the measures $X_{t_N},\nu_{t_N}$ over $\mathfrak{S}_{n_N}$ the same way as in \Cref{sec: Intro}. Then, we have 
$$\lim_{N\rightarrow\infty}d_{TV}(X_{t_N},\nu_{t_N})\cdot(\log n_N)^{5/2}=0.$$
\end{thm}

For the rest of the proof, for notational ease, we will drop off the subscript and simply write $k,n$ for $k_N,n_N$. Now, let us recall necessary background from representation theory.

\subsection{Nonabelian Fourier transform and related notations}

Given a finite group $G$, we let $\widehat G$ denote the set of irreducible representations. We define convolution $f_1,f_2:G\rightarrow\C$ to be
$$(f_1*f_2)(z):=\sum_{x\in G}f_1(x)f_2(x^{-1}z).$$
For $\rho\in \widehat G$, let $d_\rho$ denote the dimension of $\rho$. The nonabelian Fourier transform for a function $f:G\rightarrow\C$ is the map $\widehat f$ given by
$$\widehat f(\rho):=\sum_{g\in G}f(g)\rho(g)\in\Mat_{d_\rho}(\C),\quad\forall \rho\in\widehat G.$$
Then, for all $f_1,f_2:G\rightarrow\C$ and $\rho\in\widehat G$, we have
\begin{equation}\label{eq: convolution to multiplication}
\widehat{f_1*f_2}(\rho)=\widehat f_1(\rho)\cdot\widehat f_2(\rho).
\end{equation}
In this paper, we always take $G=\mathfrak{S}_n$. Following the standard construction of Specht modules, we can identify elements in $\widehat{\mathfrak{S}}_n$ with (positive integer) partitions of $n$. From now on, we use $\l\vdash n$ to denote that $\l$ is a partition of $n$. Given a partition $\l=(\l_1,\ldots,\l_j)$ with $\l_1\ge\cdots\ge\l_j$, we let $\l'$ denote the conjugate partition and $\l^*=(\l_2,\ldots,\l_j)$ be the partition of $n-\l_1$ obtained by truncating $\l$. Denote by $l(\lambda)$ the number of parts of $\lambda$, and
$$
\lambda^\bullet:=\bigl((\lambda')^*\bigr)'\,,
$$
which is obtained from $\lambda$ by deleting its first column. 

Let
\begin{equation}\label{eq: long first row and column partitions}
\mathcal L \ :=\ \Bigl\{\lambda\vdash n:\ \lambda_1\ge n-\log n\ \ \text{or}\ \ l(\l)\ge n-\log n\Bigr\}
\end{equation}
be the set of partitions with long first row or long first column. It is clear that $\mathcal L=\mathcal L_1\bigsqcup \mathcal L_2$, where
$$\mathcal L_1:=\ \Bigl\{\lambda\vdash n:\ \lambda_1\ge n-\log n\Bigr\},\quad\mathcal L_2:=\ \Bigl\{\lambda\vdash n:\ l(\l)\ge n-\log n\Bigr\}.$$
Moreover, under the conjugation map, the partitions in $\mathcal{L}_1$ have one-to-one correspondence with the partitions in $\mathcal{L}_2$.  

Probability measures on $\mathfrak{S}_n$ will be viewed as real-value functions. In our applications, we will furthermore restrict attention to functions $f$ which are class functions, i.e., functions that are constants on conjugacy classes. 
\begin{lemma}[Schur's lemma]
Suppose $f:\mathfrak{S}_n\rightarrow \C$ is a class function. Then, we have
$$\widehat f(\lambda)=\frac{\sum_{g\in G}f(g)\chi_\lambda(g)}{d_\lambda}\Id_{d_\lambda},$$
where $\chi_\lambda$ is the character corresponding to $\lambda\vdash n$.
\end{lemma}

In particular, let $P_n$ be the probability measure over $\mathfrak{S}_n$ as in \eqref{eq: measure of k-cycle}. Then, its Fourier transform is given by
\begin{equation}\label{eq: Fourier transform of Pn}
\widehat P_n(\lambda)=\frac{\chi_\lambda(\tau_k)}{d_\lambda}\Id_{d_\lambda},\quad \lambda\vdash n.
\end{equation}
where $\tau_k$ is an arbitrary $k$-cycle. Following \cite{nestoridi2022limit}, we will also use the abbreviation $s_\lambda(k)$ for $\frac{\chi_\lambda(\tau_k)}{d_\lambda}$. In particular, by the hook length formula, we have $d_\l=d_{\l'}$. Furthermore, the natural isomorphism $S^{\lambda'}\cong S^\lambda\otimes\sgn$ leads to the fact that $\chi_{\l'}(k)=\chi_\l(k)\cdot(-1)^{k-1}$, and therefore
$$s_{\l'}(k)=s_\l(k)\cdot(-1)^{k-1}.$$

\subsection{Character estimates for $P_n^{*t}$}

Following \eqref{eq: convolution to multiplication} and \eqref{eq: Fourier transform of Pn}, it is clear that
$$\widehat {P_n^{*t}}(\l)=s_\l(k)^t\Id_{d_\l},\quad\l\vdash n.$$
In the following, we provide quantitative bounds that will be useful later. Our range of $t$ will always follow \Cref{thm: X and nu}.

\begin{lemma}[{\cite[Lemma 3.3]{jain2024hitting}}]\label{lem: dimension bound}
Let $\l\vdash n$, and $r:=n-\l_1$. Then, we have
\begin{enumerate}
\item $d_\lambda\le\binom{n}{r}d_{\lambda^*}\le\binom{n}{r}\cdot \sqrt{r!}\le n^r/\sqrt{r!}.$
\item When $r\le \log n$, we have
$$d_\l=\binom{n}{r}\cdot d_{\l^*}\left(1-\frac{r}{n}+O(n^{-2+o(1)})\right).$$
\end{enumerate}
\end{lemma}

\begin{prop}[{\cite[Proposition 2.6]{shen2026k}}]\label{prop: small r}
Suppose $r:=n-\l_1\in[1, \log n]$, and $2\le k\le o(n/(\log n)^4)$. Moreover, suppose that $t':=t-\frac{n\log n}{k}$ satisfies $|t'|\le \frac{n\log\log\log n}{2k}$. Then, we have
$$s_\l(k)^t=\left(1-O\left(\frac{r(k+r)\log n}{n}\right)\right)\exp\left(-\frac{rkt}{n}\right).$$
\end{prop}

\begin{proof}
Following \cite[Proposition 2.6]{shen2026k}, we already have
$$s_\l(k)^t=\exp\left(-\frac{rkt}{n}-\frac{r(k+r)\log n}{2n}(1+o(1)+O(1/k))\right).$$
Now, since $k=o(n/(\log n)^4)$ and $r\le \log n$, we have $\frac{r(k+r)\log n}{2n}(1+o(1)+O(1/k))=o(1)$, so that
$$s_\l(k)^t=\left(1-O\left(\frac{r(k+r)\log n}{n}\right)\right)\exp\left(-\frac{rkt}{n}\right).$$
\end{proof}

The following bound generalizes \cite[Lemma 3.4]{jain2024hitting}. It shows that the sum $\sum_{\lambda\vdash n}d_\l^2|s_\l(k)|^{2t}$ is mainly contributed by partitions in $\mathcal{L}$. 

\begin{thm}[{\cite[Theorem 2.15]{shen2026k}}]\label{thm: removal}
Suppose that $2\le k\le o(n/(\log n)^4)$, and that $t':=t-\frac{n\log n}{k}$ satisfies $|t'|\le \frac{n\log\log\log n}{2k}$. Then, we have
$$\sum_{\substack{\lambda\vdash n\\\lambda\notin \mc{L}}}d_\l^2|s_\l(k)|^{2t}=n^{-\omega(1)}.$$
\end{thm}

Let us turn to character estimates for $\nu_t$. Fix integers $n\ge 1$ and $0\le M\le n/3$. In this case, for $\l\vdash n$, at most one of the following holds:
\begin{enumerate}
\item $\l_1\ge n-M$.
\item $l(\l)\ge n-M$.
\end{enumerate}

Define a probability measure $\xi_M$ (resp. $\xi_M^c$) on $\mathfrak{A}_n$ as follows: choose a subset $S\subset[n]$ uniformly among all subsets of size $M$, then sample a permutation $\pi$ uniformly from $\A_{[n]\setminus S}$ (resp. $\A_{[n]\setminus S}^c$ ), and extend it to an element of $\mathfrak{A}_n$ (resp. $\mathfrak{A}_n^c$) by fixing every point of $S$. The following lemma gives expression of the Fourier transform of $\xi_M$ and $\xi_M^c$.

\begin{thm}[Fourier transform of $\xi_M$ and $\xi_M^c$]
\label{thm :zetaM_Fourier_An}
We have
$$
\widehat{\xi_M}(\lambda)=\frac{K_{\lambda,\mu}+K_{\lambda',\mu}}{d_\lambda}\,\Id_{d_\lambda}, \quad\widehat{\xi_M^c}(\lambda)=\frac{K_{\lambda,\mu}-K_{\lambda',\mu}}{d_\lambda}\,\Id_{d_\lambda},
$$
where
$$
\mu=\mu_M:=(n-M,1,1,\dots,1)\vdash n,
$$
and $K_{\alpha,\mu}$ denotes the Kostka number.

Moreover, if $\l_1<n-M$ and $l(\l)<n-M$, then $\widehat{\xi_M}(\lambda)=\widehat{\xi_M^c}(\lambda)=0$. If $\l_1\ge n-M$, we have
$$\widehat{\xi_M}(\lambda)=\binom{M}{\,n-\lambda_1\,}\frac{
d_{\lambda^*}}{d_\lambda}\,\Id_{d_\lambda},\quad\widehat{\xi_M^c}(\lambda)=\binom{M}{\,n-\lambda_1\,}\frac{
d_{\lambda^*}}{d_\lambda}\,\Id_{d_\lambda}.$$
If $l(\l)\ge n-M$, we have
$$\widehat{\xi_M}(\lambda)=\binom{M}{\,n-l(\lambda)\,}\frac{d_{\lambda^\bullet}}{d_\lambda}\,\Id_{d_\lambda},\quad\widehat{\xi_M^c}(\lambda)=-\binom{M}{\,n-l(\lambda)\,}\frac{d_{\lambda^\bullet}
}{d_\lambda}\,\Id_{d_\lambda}.$$
\end{thm}

\begin{proof}
This is \cite[Theorem 2.16]{shen2026k}.
\end{proof}

We already have the necessary preparation to prove \Cref{thm: restate of X and nu}. 

\begin{proof}[Proof of \Cref{thm: restate of X and nu}]
In order to simplify computations, we consider the following truncated version $\mu_{n,k}^t$ of $\nu_{n,k}^t$, which is defined in a similar fashion, except that the distribution of $M_t$ (i.e., the size of the random subset $S_t$) is given by
$$\mathbf{P}[M_t=x]=\mathbf{P}[\Pois(\gamma_{n,t})=x]/\mathbf{P}[\Pois(\gamma_{n,t})\le\log n],\quad x\le \log n.$$
Here, $\gamma_{n,t}:=\exp(-kt'/n)$ follows from \Cref{defi: construction of nu}. Due to our assumption on $t'$, we have $\gamma_{n,t}\le\log\log n$, and therefore
\begin{equation}\label{eq: Poisson truncate}
\mathbf{P}[\Pois(\gamma_{n,t})\ge \log n]\le \mathbf{P}[\Pois(\log\log n)\ge \log n]=n^{-\omega(1)}.
\end{equation}
Thus, it follows from the natural decoupling that $d_{TV}(\nu_{n,k}^t,\mu_{n,k}^t)=n^{-\omega(1)}$. Thus, it suffices to prove the same bound for $d_{TV}(P_n^{*t},\mu_{n,k}^t)$. Let $f_1$ be the probability density function of $P_n^{*t}$, and $f_2$ be the probability density function of $\mu_{n,k}^t$. Using the Cauchy-Schwarz inequality followed by the Plancherel formula, we have that (the dagger superscript refers to the conjugate transpose matrix)
\begin{align}
\begin{split}
4\cdot d_{TV}(P_n^{*t},\mu_{n,k}^t)^2&\le n!\sum_{\sigma\in\S_n}(f_1(\sigma)-f_2(\sigma))^2\\
&=\sum_{\l\vdash n}d_\l\Tr\left((\widehat f_1(\l)-\widehat f_2(\l))(\widehat f_1(\l)-\widehat f_2(\l))^\dag\right)\\
&=\sum_{\l\notin\L}d_\l^2|s_\l(k)|^{2t}+\sum_{\l\in\L}d_\l\Tr\left((\widehat f_1(\l)-\widehat f_2(\l))(\widehat f_1(\l)-\widehat f_2(\l))^\dag\right)\\
&=n^{-\omega(1)}+\sum_{\l\in\L}d_\l\Tr\left((\widehat f_1(\l)-\widehat f_2(\l))(\widehat f_1(\l)-\widehat f_2(\l))^\dag\right).\\
\end{split}
\end{align}
Here, the third line follows from \Cref{thm :zetaM_Fourier_An} since $\widehat f_2(\l)=0$ when $\l\notin\L$, and the final line follows from \Cref{thm: removal}. Moreover, due to the symmetry under conjugation, the partitions in $\mathcal L_1$ and $\mathcal L_2$ contribute equally to the whole sum:
the sign change $s_{\lambda'}(k)^t=(-1)^{t(k-1)}s_\lambda(k)^t$ is matched
by the corresponding sign change in the Fourier coefficient of the parity-supported comparison measure. Hence, we only need to prove
\begin{equation*}
\sum_{\l\in\L_1}d_\l\Tr\left((\widehat f_1(\l)-\widehat f_2(\l))(\widehat f_1(\l)-\widehat f_2(\l))^\dag\right)=o\left((\log n)^{-5}\right).
\end{equation*}
Now, suppose that $\l\in\mathcal{L}_1$, and denote $r:=n-\l_1$. On the one hand, by \Cref{prop: small r}, we have
\begin{equation}
\widehat f_1(\l)=s_\l(k)^t\Id_{d_\l}=\left(1-O\left(\frac{r(k+r)\log n}{n}\right)\right)\exp\left(-\frac{rkt}{n}\right)\Id_{d_\l}.
\end{equation}
On the other hand, notice that
$$\mu_{n,k}^t\sim\begin{cases}
\sum_{l=0}^{\log n}\xi_l^c\cdot\frac{\mathbf{P}[\Pois(\gamma_{n,t})=l]}{\mathbf{P}[\Pois(\gamma_{n,t})\le\log n]} & t \text{ odd},k\text{ even}\\
\sum_{l=0}^{\log n}\xi_l\cdot\frac{\mathbf{P}[\Pois(\gamma_{n,t})=l]}{\mathbf{P}[\Pois(\gamma_{n,t})\le\log n]} & \text{else}
\end{cases},$$
where $\xi_l,\xi_l^c$ are the distributions in the statement of \Cref{thm :zetaM_Fourier_An}. Therefore, it follows from \Cref{thm :zetaM_Fourier_An} and \Cref{lem: dimension bound} that
\begin{align}
\begin{split}
\widehat f_2(\l)&=\mathbf{P}[\Pois(\gamma_{n,t})\le\log n]^{-1}\cdot\sum_{l=0}^{\log n}\frac{d_{\lambda^*}\binom{l}{r}}{d_\lambda}\mathbf{P}[\Pois(\gamma_{n,t})=l]\Id_{d_\l}\\
&=\left(1-\frac rn+O(n^{-2+o(1)})\right)\sum_{l=0}^{\log n}\binom{l}{r}\cdot r!n^{-r}\cdot\frac{\gamma_{n,t}^l \exp(-\gamma_{n,t})}{l !}\Id_{d_\l}\\
&=\left(1-\frac rn+O(n^{-2+o(1)})\right)\sum_{l=r}^{\log n}n^{-r}\cdot\frac{\gamma_{n,t}^l \exp(-\gamma_{n,t})}{(l-r)!}\Id_{d_\l}\\
&=\left(1-\frac rn+O(n^{-2+o(1)})\right)n^{-r}\cdot\gamma_{n,t}^r\cdot\mathbf{P}[\Pois(\gamma_{n,t})\le \log n-r]\Id_{d_\l}\\
&=\left(1-\frac rn+O(n^{-2+o(1)})\right)\exp\left(-rkt/n\right)\mathbf{P}[\Pois(\gamma_{n,t})\le\log n-r]\Id_{d_\l}.
\end{split}
\end{align}
Here, we use the fact from \eqref{eq: Poisson truncate} that $\mathbf{P}[\Pois(\gamma_{n,t})\le\log n]=1-n^{-\omega(1)}$. Now, suppose
$$\widehat f_1(\l)-\widehat f_2(\l)=\alpha_\l\Id_{d_\l},\quad \l\in\L_1,$$
so that 
$$\sum_{\l\in\L_1}d_\l\Tr\left((\widehat f_1(\l)-\widehat f_2(\l))(\widehat f_1(\l)-\widehat f_2(\l))^\dag\right)=\sum_{\l\in\L_1}d_{\l}^2|\alpha_{\l}|^2.$$
We break the above sum into two parts, depending on whether $r\ge\frac12\log n$ or not.
\begin{enumerate}
\item Suppose $r\ge \frac 12\log n$. In this case, we can use the rough bound $\widehat f_1(\l)-\widehat f_2(\l)=\alpha_\l\Id_{d_\l}$ where $|\alpha_\l|\le 3\exp(-rkt/n)$, along with the dimension bound $d_\l^2\le n^{2r}/r!$ in \Cref{lem: dimension bound} to see that
\begin{align}
\begin{split}
\sum_{n-\l_1\in[\frac 12\log n,\log n]}d_\l^2|\alpha_\l|^2&\le9\sum_{r\in[\frac 12\log n,\log n]}\sum_{\l_1=n-r}\exp(-2rkt/n)\frac{n^{2r}}{r!}\\
&\le O(\log n)\cdot\max_{r\in[\frac 12\log n,\log n]}\frac{\exp(-2rkt'/n)}{r!}\exp\left(\pi\sqrt{2r/3}\right)\\
&= n^{-\omega(1)}.
\end{split}
\end{align}
Here, in the last inequality, we used Stirling's formula and the assumption $|2kt'/n|\le\log\log\log n$.

\item Suppose $r<\frac 12\log n$. In this case, we have 
$$\mathbf{P}[\Pois(\gamma_{n,t})\le \log n-r]\ge \mathbf{P}\left[\Pois(\log\log n)\le \frac 12\log n\right]\ge1-n^{-\omega(1)}.$$ 
Therefore, we have $\widehat f_2(\l)=\left(1-\frac rn+O(n^{-2+o(1)})\right)\exp\left(-rkt/n\right)$, and
\begin{align*}
\begin{split}
\alpha_\lambda&=\exp(-rkt/n)\left(\left(1-O\left(\frac{r(k+r)\log n}{n}\right)\right)-\left(1-\frac rn+O(n^{-2+o(1)})\right)\right)\\
&=O\left(\frac{r(k+r)\log n}{n}\right)\exp(-rkt/n).
\end{split}
\end{align*}
Also, by \Cref{lem: dimension bound}, we obtain
\begin{equation}\label{eq: estimate of sum of d_l^2}
\sum_{\l_1=n-r}d_\l^2\le \binom{n}{r}^2\sum_{\mu\vdash r}d_\mu^2=\binom{n}{r}^2\cdot r!\le\frac{n^{2r}}{r!}.
\end{equation}
Thus, applying \eqref{eq: estimate of sum of d_l^2}, we have 
\begin{align}
\begin{split}
\sum_{n-\l_1<\frac 12\log n}d_\l^2|\alpha_\l|^2&\le \sum_{r<\frac 12\log n}O\left(\frac{r(k+r)\log n}{n}\right)^2\frac{n^{2r}}{r!}\exp(-2rkt/n)\\
&\le \sum_{r\ge 0}O\left(\frac{r(k+r)\log n}{n}\right)^2\frac{\exp(-2rkt'/n)}{r!}\\
&=O\left(\frac{k^2(\log n)^2}{n^2}\cdot\log n\right)\\
&=o\left((\log n)^{-5}\right).
\end{split}
\end{align}
Here, on the third line, we use the fact that $|2kt'/n|\le\log\log\log n$, which implies 
$$\sum_{r\ge 0}r^m\frac{\exp(-2rkt'/n)}{r!}\le\sum_{r\ge 0}r^m\frac{(\log\log n)^r}{r!}=O(\log n)$$
for integers $m=2,3,4$.
\end{enumerate}

Combining our discussion of $r\in[1,\frac 12\log n)$ and $r\in[\frac 12\log n,\log n]$ above, we complete the proof.
\end{proof}




